\section{From search records to a defensible coverage claim}
\label{exec13:survey-coverage}

A systematic comparison needs a rule for what one entry represents. Let
$\mathcal P$ be the set of canonical implementations, $\mathcal A$ the method
areas and $\mathcal L$ the language or commercial access paths. A source may
support a relation $(p,a,\ell)\in\mathcal P\times\mathcal A\times\mathcal L$.
This is a many-to-many relation. One implementation can supply several methods;
one method can have several implementations; and a wrapper can expose an
implementation in another language. Counting each relation as a new package
would overstate the number of independent alternatives. Collapsing all
relations to the package name would conceal unsupported methods and access
paths. The survey therefore needs both a canonical implementation index and
an operation-specific comparison.

The pcalg example illustrates this distinction without requiring a hypothetical
new package. Its PC discovery and IDA effect calculation belong to the same
implementation family, but their inputs and outputs differ. Reading the PC
interface does not establish an IDA claim. Similarly, the documented R
interface does not establish that a proposed Python process wrapper has
preserved its graph semantics. The inspected manual supports the named native
operations; the wrapper requires additional evidence
\citep{exec13_pcalg}. The canonical entry can share identity and license facts
while keeping separate evidence for each operation.

For each comparison, define the claim as a tuple
\[
 c=(p,v,a,\ell,d,q,o),
\]
where $v$ is a release or source revision, $d$ the supported data regime,
$q$ the scientific request and $o$ the output meaning. A useful source record
answers a specific question about this tuple. A manual can describe an
operation on complete iid continuous records without establishing behavior
on clustered monthly histories. A license can establish a permission subject
to conditions without establishing algorithmic correctness. An executed
roundtrip can establish endpoint preservation without establishing coverage
of a confidence interval. These evidence types are complementary; none
should replace the others in a single undifferentiated confidence score.

The comparison can use five explicit evidence states: a discovered lead;
an identified primary source; a technically inspected claim; an executed
behavioral reference; and a statistically evaluated claim in a declared law.
An entry moves between these states only for the relevant claim tuple.
Failure is retained separately: an inaccessible manual is an acquisition
failure, an unsupported input is an interface restriction, and an invalid
effect calculation is a substantive problem. Each implies a different next
action. A search result identifying a package cannot advance a numerical
claim, and a successful trial cannot establish that an undocumented edition
shares the same feature.

This structure also clarifies survey coverage. The fraction of predefined
comparison cells with a reviewed disposition measures completion of that
declared plan. It does not estimate the fraction of all software in the world
that has been found. Ranked bibliographic results have unknown inclusion
probabilities, aliases can occupy several results, and broad queries return
method papers without an implementation. Dividing ten inspected records by
the 1,648 reported matches would therefore provide no defensible software
recall estimate. The bounded R comparison in this chapter reports its search
route and five primary-inspected paths, without attaching a census percentage.

A practical extension procedure is to search exact package names and aliases,
method names, implementation languages and competing commercial descriptions;
then inspect the primary technical sources for the resulting candidates.
Reference following adds implementations that keyword choice may miss.
Backward and forward search have different date constraints: a later search
service can locate an earlier work, but a post-cutoff technical revision must
not be presented as the earlier implementation. Conflicting manuals or
source versions create a comparison question that needs resolution at a
specific version. A later edit cannot silently repair what a preserved
earlier inspection established.

The stopping rule should be chosen before interpreting omissions. A bounded
proposal pass may inspect a fixed number of query pages and one reference
generation, retain every discovered identity, and defer unresolved leads.
Those are resource limits on the search, not absence findings about a market.
For a consequential adoption choice, the stopping condition is stricter:
each proposed operation needs its technical source, deployment and rights
path, alternatives and a testable behavior contract. If an essential source
is missing, that adoption decision stays conditional while other comparisons
continue. This makes the unknown part of the landscape visible and gives
the next search a reason tied to a decision.

\section{A matched comparison includes eligible requests and justified refusal}
\label{exec13:survey-comparison}

The common contract makes it possible to define a fair comparison before
measuring a package. First choose a service regime: for example, Gaussian
iid discovery and singleton linear effects, with the same named variables,
sampling unit, units and background restrictions for every candidate. Then
separate three case families. In the first, the target is identified and a
reference is available. In the second, several causal models remain
observationally compatible and an effect set is the reference. In the third,
the request violates the service regime or lacks necessary input. A correct
response in the last two families can be an unresolved set or a typed refusal.
Forcing every case into a point-estimate error would reward unsupported
answers and penalize scientifically useful ambiguity.

Let $z_i$ indicate that case $i$ has a reference answer within the agreed
service, and let $r_{ij}$ indicate that path $j$ returned a valid answer.
Eligible availability is
\[
 A_j=\frac{\sum_i z_i r_{ij}}{\sum_i z_i}.
\]
An unavailable estimate on an eligible case remains in this denominator.
For cases outside the service, separately record the fraction receiving an
unjustified causal answer. An invocation that completed successfully but
changed the population, units or graph class belongs in the latter error
category when those changes make its answer unsupported. A library exception
and a justified scientific refusal are also different outcomes: the former
can indicate a broken supported path, whereas the latter can be the intended
response to an unsupported request.

For each successfully answered identified case, numerical error and interval
quality have their own references. A small error conditional on returning
an answer does not establish good service if difficult cases fail. Conversely,
high availability does not establish valid inference. Report the joint event
of an available interval covering the target, its conditional coverage when
available, and width, alongside failures. For an effect-set case, compare
the complete required object: missing an admissible effect and unnecessarily
adding impossible effects are different defects. If a benchmark grants
credit for preserving a convex hull, it must say so; that does not test
preservation of a disconnected set or multiplicity.

An exact planning example shows why these distinctions can change a choice.
Suppose a generated comparison contains 100 requests, with 60 inside the
shared service and 40 requiring refusal under its declared scope. These
counts are stipulated for illustration. Path A answers 54 eligible cases
correctly, fails on six, gives eight unjustified answers and correctly refuses
32. Path B answers 48 eligible cases correctly, fails on twelve and correctly
refuses all 40 ineligible cases. Eligible availability is 90\% for A and
80\% for B, while false acceptance among ineligible requests is 20\% and
zero. Their valid-answer fractions among all requests are 54\% and 48\%.
No single completion rate contains all this information
\citep{exec13_survey_comparison}.

To connect the same example to a decision, let a valid answer have expected
incremental value $v$, an unjustified answer expected loss $h$, and the
full operating costs for the 100 requests be $c_A$ and $c_B$. The assumed
net values are
\[
 V_A=54v-8h-c_A,\qquad V_B=48v-c_B.
\]
With hypothetical currency values $v=100$, $c_A=400$ and $c_B=600$, the
two paths tie at $h=100$ with net value 4,200. A is preferred only below
that harm threshold, under this deliberately simple loss model. The counts,
prices and values are project assumptions, reproduced by exact arithmetic;
they are neither vendor measurements nor observed bank returns
\citep{exec13_survey_comparison}. The example makes a practical point:
extra answered requests have to be valued together with the consequences
of answering an unsupported question.

Actual prototype comparisons should therefore retain individual case results,
not merely aggregates. The case identity carries the declared law and target;
the result records the returned object, failures and resource use. Paired
comparisons use the same cases and distinguish a numerical discrepancy from
a different target. Tuning and adapter repairs use development cases, while
a later evaluation freezes the selected procedure. If a package changes the
search space or target to obtain a result, that becomes a different procedure
to evaluate. These rules allow a favorable or adverse finding to be traced
to a concrete operation and make a new trial useful when a source version
or proposed wrapper changes.

The same protocol informs reuse cost. The full cost includes acquisition and
dependency work, conversion, analyst interpretation, failed requests and
future maintenance. A free download can still require expensive scientific
adaptation; an enterprise license can be worthwhile if its documented service
and measured workflow fit the request. Before paying for a trial, the proposal
can identify the unresolved comparison and calculate the decision threshold
at which its answer matters. The next expense then purchases information
about a specific choice. The present source survey supplies candidate paths
and testable restrictions; the proposed matched trial supplies the evidence
needed to choose among them.
\label{exec13:survey-reviewed-end}
